\section*{Spivak Capítulo 3, Problema 10}
\addcontentsline{toc}{section}{Spivak Capítulo 3, Problema 10}

El ejercicio plantea preguntas sobre la existencia de funciones que satisfacen ciertas ecuaciones punto a punto. Sea $\mathcal{F}$ el conjunto de todas las funciones reales definidas en un mismo dominio $D\subseteq\mathbb{R}$ (normalmente todos los reales o algún intervalo).

\begin{enumerate}[label=(\alph*)]
\item ¿Para qué funciones $f\in\mathcal{F}$ existe otra $g\in\mathcal{F}$ tal que
$$f(t)=\bigl(g(t)\bigr)^2\qquad\text{para todo }t?$$

\textbf{Respuesta.} Una función $f$ admite una "raíz cuadrada" $g$ punto a punto si y solo si $f(t)\ge0$ para todo $t$; en ese caso podemos tomar
$$g(t)=\sqrt{f(t)}$$
(o la negativa $g(t)=-\sqrt{f(t)}$, o elegir signo distinto en cada $t$ si se desea). Si hay algún $t$ con $f(t)<0$ no existe $g$ real satisfaciendo la ecuación en ese $t$.

\item ¿Para qué funciones $f$ existe $g$ tal que $f(t)=1/g(t)$ para todo $t$?

\textbf{Respuesta.} Esto equivale a pedir una función recíproca. Se requiere que $f(t)\neq0$ para todo $t$; entonces basta definir $g(t)=1/f(t)$. Si $f$ anula en algún punto no hay función $g$ real que cumpla la igualdad allí.

\item (marcado *) ¿Para qué funciones $b,c$ podemos encontrar una función $x$ tal que
$$(x(t))^2+b(t)x(t)+c(t)=0\qquad\text{para todo }t?$$

\textbf{Respuesta.} Para cada $t$ esto es una ecuación cuadrática en la variable $x(t)$. Tiene solución real si y sólo si el discriminante es no negativo:
$$b(t)^2-4c(t)\ge0\quad\text{para todo }t.$$  En ese caso se pueden elegir puntualmente
$$x(t)=\frac{-b(t)\pm\sqrt{b(t)^2-4c(t)}}2.$$  Si en algún punto el discriminante es negativo no existe solución real completa.

\item (marcado *) ¿Qué condiciones deben satisfacer las funciones $a$ y $b$ si ha de existir una función $x$ tal que
$$a(t)x(t)+b(t)=0\qquad\text{para todo }t?$$
¿Cuántas funciones $x$ de éstas existirán?

\textbf{Respuesta.} Para que exista $x$ basta que $a(t)\neq0$ para todos los $t$ (si $a(t)=0$ y $b(t)\neq0$ no se puede resolver, si ambos son cero hay infinitas soluciones locales). Cuando $a(t)\neq0$ definimos
$$x(t)=-\frac{b(t)}{a(t)}$$
única en cada punto. En el caso más general, si $A=\{t:a(t)=0\}$, entonces
\begin{itemize}
\item si $b(t)\neq0$ para algún $t\in A$, no existe $x$ global;
\item si $b(t)=0$ para todos $t\in A$, entonces hay infinitas soluciones: en los puntos de $A$ podemos asignar cualquier valor a $x(t)$, mientras que fuera de $A$ está forzada por la fórmula anterior.
\end{itemize}

\end{enumerate}

Estas observaciones se obtienen simplemente pensando en el caso numérico (sustituir "función" por "número"), tal y como sugiere el enunciado.